\subsection{Test di sistema}
\label{subsec:test_sistema}

\begin{tabularx}{\textwidth}{|c|X|c|}
    \hline
    \textbf{ID} & \textbf{Descrizione} & \textbf{Stato}\\
    \hline
    \label{ts:ts17} TS.17 & Verificare che l'utente possa visualizzare nome e data dell'ultimo aggiornamento del dataset caricato & \\
    \hline
    \label{ts:ts18} TS.18 & Verificare che l'utente possa rinonimare un dataset &  \\
    \hline
    \label{ts:ts19} TS.19 & Verificare che l'utente visualizzi un messaggio di errore se il nome per il dataset inserito non è valido, cioè se è: %17.04 
                            \begin{itemize}
                                \item vuoto;
                                \item composto da soli spazi;
                                \item già utilizzato per un altro dataset.
                            \end{itemize}  
                            & \\ 
    \hline
    \label{ts:ts20} TS.20 & Verificare che l'utente possa caricare un dataset sovrascrivendo il dataset attualmente caricato & \\
    \hline
    \label{ts:ts21} TS.21 & Verificare che il sistema, caricato un test, gli associ correttamente nome e versione del dataset utilizzato & \\ 
    \hline
    \label{ts:ts22} TS.22 & Verificare che il sistema gestisca correttamente la paginazione dei risultati del test caricato & \\ 
    \hline
    \label{ts:ts23} TS.23 & Verificare che l'utente visualizzi correttamente i risultati del test, l'utente deve visualizzare:
                            \begin{itemize}
                                \item i risultati del test per ogni domanda del dataset testato, quindi la risposta ottenuta e il grado di similarità;
                                \item la media e la deviazione standard dei risultati del test;
                                \item i risultati del test raggruppati nei seguenti range: [0, 0.2], [0.2, 0.4], [0.4, 0.6], [0.6, 0.8], [0.8, 1];
                                \item la percentuale di risposte corrette ottenute durante il test.
                            \end{itemize}  & \\ 
    \hline
    \label{ts:ts24} TS.24 & Verificare che l'utente possa visualizzare una lista scorrevole di tutti gli LLM possibili per il testing & \\ 
    \hline
    \label{ts:ts25} TS.25 & Verificare che l'utente possa visualizzare per ogni LLM salvato:
                            \begin{itemize}
                            \item nome e data di ultimo aggiornamento;
                            \item elementi di configurazione per la chiamata HTTP all'LLM.
                            \end{itemize} & \\ 
    \hline
    \label{ts:ts26} TS.26 & Verificare che l'utente possa selezionare l'LLM da testare prima di iniziare il test & \\ %RFF-1.03
    \hline
    \label{ts:ts27} TS.27 & Verificare che l'utente visualizzi un messaggio di errore nel caso ci siano stati problemi nell'ottenimento delle informazioni degli LLM & \\ %RFF-1.04
    \hline
    \label{ts:ts28} TS.28 & Verificare che l'utente possa eliminare un LLM salvato & \\ %RFF-2.03
    \hline
    \label{ts:ts29} TS.29 & Verificare che l'utente, prima dell'eliminazione dell'LLM salvato, visualizzi:
                            \begin{itemize}
                                \item elenco dei test realizzati sull'LLM che verranno eliminati di conseguenza;
                                \item richiesta di conferma dell'eliminazione.
                            \end{itemize} & \\ %RFF-2.04 2.05
    \hline
    \label{ts:ts30} TS.30 & Verificare che l'utente possa modificare un LLM salvato e quindi possa:
                            \begin{itemize}
                                \item rinonimare l'LLM;
                                \item modificare elementi di configurazione per la chiamata HTTP all'LLM;
                                \item aggiungere elementi di configurazione per la chiamata HTTP all'LLM. 
                            \end{itemize} & \\ %RFF-3.08 RFF4.01 <->4.10
    \hline
    \label{ts:ts31} TS.31 & Verificare che l'utente possa aggiungere un nuovo LLM testabile inserendo:
                            \begin{itemize}
                                \item nome del nuovo LLM;
                                \item elementi di configurazione per la chiamata HTTP all'LLM.
                            \end{itemize} & \\ %RFF-2.06 5.01 <->5.08
    \hline
    \label{ts:ts32} TS.32 & Verificare che l'utente visualizzi un messaggio di errore nel caso gli elementi 
    di configurazione per la chiamata HTTP all'LLM siano in un formato non valido & \\ %RFF-4.11 <->4.13 5.09 <->5.11
\hline
\end{tabularx}


\subsubsection{Tracciamento test di sistema}

\begin{table}[H]
    \centering
    \resizebox{\textwidth}{!}{
    \begin{tabular}{| l | c |}
    \hline
        \textbf{ID} & 
        \textbf{Requisito}\\
    \hline
        TS.17 & RFO-17.02 RFO-17.01 \\
    \hline
        TS.18 & RFO-17.03 \\
    \hline
        TS.19 & RFO-17.04 \\
    \hline
        TS.20 & RFO-17.05 \\
    \hline
        TS.21 & RFO-18.01, RFO-18.02 \\
    \hline
        TS.22 & RFO-18.03, RFO-18.04, RFO-18.05, RFO-18.06 \\
    \hline
        TS.23 & RFO-19.01, RFO-19.02, RFO-19.03, RFO-19.04, RFO-19.05, RFO-19.06, RFO-19.07\\
    \hline
        TS.24 & RFF-1.01\\
    \hline
        TS.25 & RFF-1.02, RFF-2.01, RFF-2.02, RFF-3.01, RFF-3.02, RFF-3.03, RFF-3.04, RFF-3.05, RFF-3.06, RFF-3.07\\
    \hline
        TS.26 & RFF-1.03\\
    \hline
        TS.27 & RFF-1.04\\
    \hline
        TS.28 & RFF-2.03\\
    \hline 
        TS.29 & RFF-2.04, RFF-2.05\\
    \hline 
        TS.30 & RFF-3.08, RFF-4.01, RFF-4.02, RFF-4.03, RFF-4.04, RFF-4.05, RFF-4.06, RFF-4.07, RFF-4.08, RFF-4.09, RFF-4.10\\
    \hline 
        TS.31 & RFF-2.06, RFF-5.01, RFF-5.02, RFF-5.03, RFF-5.04, RFF-5.05, RFF-5.06, RFF-5.07, RFF-5.08\\
    \hline
        TS.32 & RFF-4.11, RFF-4.12, RFF-4.13, RFF-5.09, RFF-5.10, RFF-5.11\\
    \hline
\end{tabular}}
\caption{Tracciamento test di sistema}
\label{tab:tracciamento_test_sistema} 
\end{table}
